% !PRINTED 83
\weberpage{83}
%<<CONT
If, then, $J$ were not divisible by $2^{\lambda-2}$, all the $\beta_2$ would have to be divisible by 2, which again is contrary to requirement IV.

e) \gesperrt{If $q$ is a prime that divides $m$ simply, then $q-1$ must be divisible by at least one of the primes dividing $J$.}

For if $c_1=q-1$ were relatively prime to $J$, all the $\beta_1$ would have to be divisible by $q-1$, in contradiction with IV. The further conditions for $m$ will emerge later.

\gesperrt{Determination of the Group $\mathfrak{B}$:}

The group $\mathfrak{B}$ is determined when its basis $g_1, g_2, \dots g_\nu$ is given. We shall determine the elements of the basis by their indices, and accordingly introduce the following notation. Let
\weq{7}{
\begin{aligned}
&\gamma_{1,1}, \gamma_{1,2}, \dots \gamma_{1,\mu} \quad \text{the indices of } g_1\\
&\gamma_{2,1}, \gamma_{2,2}, \dots \gamma_{2,\mu} \quad \text{\textquotedbl{}} \quad \text{\textquotedbl{}} \quad \text{\textquotedbl{}} \quad g_2\\
&\dots\dots\dots\dots\dots\dots\dots\dots\dots\dots\dots\dots\dots\dots\\
&\gamma_{\nu,1}, \gamma_{\nu,2}, \dots \gamma_{\nu,\mu} \quad \text{\textquotedbl{}} \quad \text{\textquotedbl{}} \quad \text{\textquotedbl{}} \quad g_\nu.
\end{aligned}
}
The indices $\beta_1, \beta_2, \dots \beta_\mu$ of an arbitrary element $b$ are then obtained, by (3), in the form:
\weq{8}{
\begin{aligned}
\beta_1 &\equiv \gamma_{1,1}x_1+\gamma_{2,1}x_2+\dots+\gamma_{\nu,1}x_\nu\ (\mathrm{mod}\ c_1)\\
\beta_2 &\equiv \gamma_{1,2}x_1+\gamma_{2,2}x_2+\dots+\gamma_{\nu,2}x_\nu\ (\mathrm{mod}\ c_2)\\
&\dots\dots\dots\dots\dots\dots\dots\dots\dots\dots\dots\dots\dots\dots\\
\beta_\mu &\equiv \gamma_{1,\mu}x_1+\gamma_{2,\mu}x_2+\dots+\gamma_{\nu,\mu}x_\nu\ (\mathrm{mod}\ c_\mu),
\end{aligned}
}
and we now ask for the necessary and sufficient conditions on the numbers $\gamma_{h,k}$ in order that $g_1, g_2 \dots g_\nu$ be the basis of a group $\mathfrak{B}$ with the required properties. By 1., it is first of all necessary for this that:

\gesperrt{2. The numbers $\gamma_{h,k}$ must be so constituted that the congruences}
\weq{9}{
\begin{aligned}
\gamma_{1,1}x_1+\gamma_{2,1}x_2+\dots+\gamma_{\nu,1}x_\nu &\equiv 0\ (\mathrm{mod}\ c_1)\\
\gamma_{1,2}x_1+\gamma_{2,2}x_2+\dots+\gamma_{\nu,2}x_\nu &\equiv 0\ (\mathrm{mod}\ c_2)\\
&\dots\dots\dots\dots\dots\dots\dots\dots\dots\dots\dots\dots\\
\gamma_{1,\mu}x_1+\gamma_{2,\mu}x_2+\dots+\gamma_{\nu,\mu}x_\nu &\equiv 0\ (\mathrm{mod}\ c_\mu)
\end{aligned}
}
\gesperrt{are satisfied then and only then. when at the same time the congruences}
% [sic: Weber prints a full stop before the lower-case ``wenn''; a comma is grammatically required.]
\weq{10}{x_1\equiv 0\ (\mathrm{mod}\ i_1),\quad x_2\equiv 0\ (\mathrm{mod}\ i_2),\quad \dots\ x_\nu\equiv 0\ (\mathrm{mod}\ i_\nu)}
\gesperrt{hold.}
